\documentclass[11pt]{article}
\usepackage[a4paper,margin=25mm]{geometry}
\usepackage{amsmath,amssymb,amsthm,mathrsfs,xfrac,xspace,hyperref}
\newcommand{\ie}{i.e.,\xspace}
\newcommand{\eg}{e.g.,\xspace}
\newcommand{\etc}{etc.\xspace}
\usepackage{mathrsfs}
\let\mathcal\mathscr
\multlinegap0pt
\newcounter{case}
\newcounter{subcase}

\newcommand{\vol}{\mathrm{vol}}

\newcommand{\Changel}{\mathcode`l="7160}
\newcommand{\Changelback}{\mathcode`l="716C}

\newcommand{\NoChangel}[1]{%
\expandafter\let\csname old\string#1\endcsname=#1
\let#1=\relax
\newcommand{#1}{\mathcode`l="716C\csname old\string#1\endcsname\mathcode`l="7160 }%
}

\NoChangel{\vol}
\NoChangel{\log}
\NoChangel{\ln}
\NoChangel{\lim}
\NoChangel{\limsup}
\NoChangel{\liminf}
\NoChangel{\varlimsup}
\NoChangel{\varliminf}
\NoChangel{\varinjlim}
\NoChangel{\varprojlim}

\let\epsilon\varepsilon
\newcommand{\Psfrac}[2]{(\sfrac{#1}{#2})}
\newcommand{\sep}{\,:\,}
\DeclareMathOperator{\Aut}{Aut}
\DeclareMathOperator{\Id}{Id}
\DeclareMathOperator{\proj}{proj}
\newcommand{\cont}{\mathrm{cont}}
\theoremstyle{definition}
\newtheorem{definition}{Definition}[section]
\newtheorem{remark}[definition]{Remark}

\theoremstyle{plain}
\newtheorem{theorem}{Theorem}
\newtheorem{proposition}[definition]{Proposition}
\newtheorem{corollary}[definition]{Corollary}
\newtheorem{lemma}[definition]{Lemma}

\newcommand\CC{\mathbb{C}}
\newcommand\NN{\mathbb{N}}
\newcommand\QQ{\mathbb{Q}}
\newcommand\RR{\mathbb{R}}
\newcommand\TT{\mathbb{T}}
\newcommand\ZZ{\mathbb{Z}}

\newcommand\cB{\mathcal{B}}
\newcommand\cC{\mathcal{C}}
\newcommand{\cD}{\mathcal D}

\newcommand{\be}{\begin{equation*}}
\newcommand{\ee}{\end{equation*}}

\begin{document}
\begin{center}{\Large Stable item 1917: Disjointness criterion}\end{center}
\paragraph{Source and attribution.}
Giovanni Forni and Adam Kanigowski, \emph{Multiple mixing and disjointness for time changes of bounded-type Heisenberg nilflows},
Journal de l'\'Ecole polytechnique --- Math\'ematiques 7 (2020), 63--91,
\href{https://doi.org/10.5802/jep.111}{DOI 10.5802/jep.111}, publisher version.
Author text licensed \href{https://creativecommons.org/licenses/by/4.0/}{CC BY 4.0}.
This compilation reproduces the authors' native TeX, with an independent layout wrapper, extracted bibliography, and explanatory source boundaries. Mathematical text is unchanged.
\paragraph{Problem boundary (editorial).}
Prove the authors' Proposition 4.1 reproduced below. Only this disjointness criterion is the selected problem; applications to nilflows and other article results are outside this item.
\paragraph{Difficulty (editorial): H2.}
A non-product ergodic joining is separated by a nonzero time displacement. Measurable selection and Lusin uniform continuity lift a positive-measure family of near-identity sheared pairs into the joining; uniform ergodic averages and compact control of the common time lag force contradictory invariance estimates. The complete argument occurs in Proposition 4.1 and its proof, beyond routine qualification-exam applications of ergodic theorems.
\paragraph{Author text: necessary definitions, original Section 2.1.}
We refer the reader to \cite{Gla} for basic theory of joinings.
Let $(\phi_t):(X,\cB,\mu)\to (X,\cB,\mu)$ and $(\psi_t):(Y,\cC,\nu)\to (Y,\cC,\nu)$ be two ergodic flows. A \emph{joining} of $(\phi_t)$ and $(\psi_t)$ is any $(\phi_t\times \psi_t)$ invariant measure such that $\rho(X\times B)=\mu(X)\nu(B)$ and $\rho(C\times Y)=\mu(C)\nu(Y)$. The set of joinings of $(\phi_t)$ and $(\psi_t)$ is denoted by $J((\phi_t),(\psi_t))$. Notice that $\mu\otimes \nu \in J((\phi_t),(\psi_t))$. We say that $(\phi_t)$ and $(\psi_t)$ are \emph{disjoint} (denoting $(\phi_t)\perp (\psi_t)$) if $J((\phi_t),(\psi_t))=\{\mu\otimes \nu\}$.

\setcounter{equation}{1}\setcounter{section}{3}
\section{Disjointness criterion}\label{sec:disj}
Let $(\phi_t):(X,\cB,\mu,d_1)\to (X,\cB,\mu,d_1)$ and $(\psi_t):(Y,\cC,\nu,d_2)\to (Y,\cC,\nu,d_2)$ be two weakly mixing flows (and $X,Y$ are $\sigma$-compact). In this section we prove the following proposition:

\begin{proposition}[Disjointness criterion]\label{disjoint.flows}
Let $P'\subset \RR$ be a compact set and let $v'\neq 0$. Fix $1\geq c>0$. Assume there exists $(A_k)\subset \Aut(X_k,\cB|_{X_k},\mu|_{X_k})$, such that $\mu(X_k)\to \mu(X)$,
$A_k\to \Id$ uniformly.
Assume moreover that for every $\epsilon>0$ and $N\in \NN$
there exist $(E_k=E_k(\epsilon))\subset \cB$ with $\mu(E_k)\geq c\mu(X)$, $0<\kappa=\kappa(\epsilon)<\epsilon$, $\delta=\delta(\epsilon,N)>0$, a set $Z=Z(\epsilon,N)\subset Y$ with $\nu(Z)\geq (1-\epsilon)\nu(Y)$, such that for all $y,y'\in Z$ satisfying $d_2(y,y')<\delta$, every $k$ such that $d_1(A_k,\Id)<\delta$ and every $x\in E_k$, $x':=A_{k}x$, there are $M\geq N$, $L\geq 1$, $\sfrac{L}{M}\geq \kappa$ and $V\in P'$, $v\in\{-v',v'\}$,
for which the following holds:\vspace*{-5pt}
\begin{multline}\label{forw}
\max(d_1(\phi_tx,\phi_{t+V+v}x'),\;d_2(\psi_ty,\psi_{t+V}y'))<\epsilon\\ \text{ for } t\in U\subset [M,M+L]
\text{ with } \lambda_\RR(U)\geq (1-\epsilon)L.
\end{multline}
Then $(\phi_t)$ and $(\psi_t)$ are disjoint.
\end{proposition}

The proof of the above proposition follows similar lines to the proof of \cite[Th.\,3]{KLU}.
We provide a proof for completeness.
\begin{proof}
Let $\rho\in J((\phi_t)_{t\in\RR},(\psi_t)_{t\in \RR})$ be an ergodic joining with $\rho\neq \mu\times \nu$. Since $(\phi_t)$ is weakly mixing it follows that, for $v\in\{-v',v'\}$, the map $\phi_v$ is ergodic and hence disjoint from $\Id$. Therefore, there exist $B_v\in \cB$ and $C_v\in \cC$ such that
\begin{equation}\label{iddisT}|\rho(\phi_{-v}(B_v)\times C_v)-\rho(B_v\times C_v)|>\eta
\end{equation}
for some $0<\eta<1$.
Let $V^1_{\epsilon}(B_v):=\{x\in X\sep d_1(x,B_v)<\epsilon\}$ and similarly $V^2_{\epsilon}(C_v):=\{y\in Y\sep d_2(y,C_v)<\epsilon\}$.
There exists $\epsilon\in(0,\sfrac{c\eta}{1000})$ such that
\[\max\left(\left|\mu(V^1_{\epsilon}(B_v))-\mu(B_v)\right|,
\left|\nu(V^2_{\epsilon}(C_v))-\nu(C_v)\right|\right) <\eta/32.\] Since $\rho$ is a joining, by the triangle inequality, for each $t\in\RR$, we have
\begin{equation}\label{pertu}
|\rho(\phi_{-t}V^1_{\epsilon}(B_v)\times V^2_\epsilon(C_v))-\rho(\phi_{-t}B_v\times C_v)|<\frac{\eta}{16}.
\end{equation}

By applying Birkhoff point-wise ergodic theorem to the joining flow $(\phi_t\times \psi_t, \rho)$ and to the characteristic functions of the sets
$\phi_{-v}B_v\times C_v$ and $\phi_{-v}V^1_{\epsilon}(B_v)\times V^2_{\epsilon}(C_v)$
for $v\in \{-v',v'\}$, it follows that there exist $N_0\in \NN$, $\kappa>0$ and a set $U_1\in \mathcal{B}\otimes \mathcal{C}$, with $\rho(U_1)>(1-\sfrac{c}{100})\rho(X\times Y)$,
such that, for every $L,M\geq N_0$ with $\sfrac{L}{M}\geq \kappa$ and $v\in \{-v',v'\}\cup\{0\}$, and for all $(x,y)\in U_1$, we have
\begin{gather}\label{eeq1}\biggl|\frac{1}{L}\int_{M}^{M+L}\chi_{\phi_{-v}B_v
\times C_v}(\phi_tx,\psi_ty)\,dt-\rho(\phi_{-v}B_v\times C_v)\biggr|<\frac{\eta}{16},
\\
\label{eeq3}\biggl|\frac{1}{L}\int_{M}^{M+L}
\chi_{\phi_{-v}V^1_{\epsilon}(B_v)
\times V^2_{\epsilon}(C_v)}(\phi_tx,\psi_ty)\,dt-\rho(\phi_{-v}V^1_{\epsilon}(B_v)\times V^2_\epsilon(C_v))\biggr|<\frac{\eta}{16}.
\end{gather}

Let $U_2:=U_1\cap (X\times Z)$, where $Z=Z(\epsilon,N_0)$ comes from our assumptions. Then $\rho(U_2)>(1-c/50)\rho(X\times Y)$.
Note also that since $X\times Y$ is $\sigma$-compact, the measure~$\rho$ is regular, and hence we can additionally assume that $U_2$ is compact.
Define $\proj:X\times Y\to X$, $\proj(x,y)=x$. Then the fibers of $\proj$ are $\sigma$-compact, and since $U_2$ is compact, the fibers of the map $\proj|_{U_2}:U_2\to \proj(U_2)\subset X$ are also $\sigma$-compact and $\proj(U_2)$ is also compact. Thus, by Kunugui's selection theorem (see \eg \cite[Th.\,4.1]{Ga-Le-Sch}), it follows that there exists a measurable (selection) $s_Y :\proj(U_2)\to X\times Y$ such that $(x,s_Y(x))\in U_2$. Note that $\mu(\proj(U_2))\geq \rho(U_2)>(1-c/50)\mu(X)$. By~Luzin's theorem there exists $X_\cont\subset \proj(U_2)$, with $\mu(X_\cont)\geq (1-c/50)\mu(X)$, such that $s_Y$ is uniformly continuous on $X_\cont$. Finally, we set
\[
\widetilde{U}:=U_2\cap (X_\cont\times Y).
\]
We have $\rho(\widetilde{U})>(1-c/10)\rho(X\times Y)$. Moreover, for the set $U_X:=\proj(\widetilde{U})$ we have
$\mu(U_X)\geq \rho(\widetilde{U})>(1-c/10)\rho(X\times Y)$.
Hence, by the definitions of sequences $(A_k)$ and $(E_k)=(E_k(\epsilon))$, it follows that there exists $k_0=k_0(\epsilon)$ such that for $k\geq k_0$,
\begin{equation}\label{kurr}
\mu(A_{-k}(U_X\cap X_k)\cap (U_X\cap X_k)\cap E_k)>0.
\end{equation}
Let $\delta=\delta(\epsilon,N_0)$ come from the assumptions of our theorem. By the uniform continuity of $s_Y:X_\cont\to Y$ it follows that there exists $0<\delta'<\delta$ such that $d_1(x_1,x_2)<\delta'$ implies $d_2(s_Y(x_1),s_Y(x_2))<\delta$ for each $x_1,x_2\in X_\cont$.
Since $A_k\to \Id$ uniformly
and $\widetilde{U}\subset X_\cont\times Y$, there exists $k_1=k_1(\epsilon)$ such that for $k\geq k_1$, $d_2(s_Y(x),s_Y(A_kx))<\delta$ for $x\in X_k\cap X_\cont$. Fix $k\geq \max(k_0,k_1+1)$ (so that $d_1(A_k,\Id)<\delta'$). Let $x\in A_{-k}(U_X\cap X_k)\cap (U_X\cap X_k)\cap E_k$. Such a point does exist in view of~\eqref{kurr}. Set $x'=A_kx$, $y=s_Y(x)$, $y'=s_Y(x')$. By definition, $(x,y),(x',y')\in \widetilde{U}$ and $d_2(y,y')<\delta$ and all other assumptions of our theorem are satisfied for $(x,y)$, $(x',y')$ (so that we obtain $M,L,V,v$ depending on $(x,y)$ and $(x',y')$ satisfying~\eqref{forw}).

We claim that
\begin{equation}\label{disjon}
\rho(\phi_{-v}(B_v)
\times C_v)>\rho(B_v\times C_v)-\frac{\eta}{2}.
\end{equation}
Indeed, in view of~\eqref{pertu}, the estimate~\eqref{disjon} follows if we can prove that
\[
\rho(\phi_{-v}(V^1_{\epsilon}(B_v))
\times V^2_{\epsilon}(C_v))>\rho(B_v\times C_v)-\frac{\eta}{4}.
\]

Using~\eqref{forw} (for $v=v'$),~\eqref{eeq1} (for $v=0$) and $(x,y)\in \widetilde{U}\subset U_1$, we obtain that
\begin{multline*}
\frac{1}{L}\int_{M}^{M+L}\chi_{V^1_{\epsilon}(B_v)
\times V^2_{\epsilon}(C_v)}(\phi_{t+V+v}x',\psi_{t+V}y')\,dt \\ > \frac{1}{L}\int_{M}^{M+L}\chi_{B_v
\times C_v}(\phi_tx,\psi_ty)\,dt-\epsilon >
\rho(B_v\times C_v)-\epsilon-\frac{\eta}{16}.
\end{multline*}
Hence to complete the proof of Claim~\eqref{disjon}, it is enough to show that
\begin{multline}\label{eq:suf}
\frac{1}{L}\int_{M}^{M+L}\chi_{V^1_{\epsilon}(B_v)
\times
V^2_\epsilon(C_v)}(\phi_{t+V+v}x,\psi_{t+V}y)\,dt<
\rho(\phi_{-v}V^1_{\epsilon}(B_v)
\times V^2_\epsilon(C_v))+\frac{\eta}{8}.
\end{multline}
Notice however that
\begin{multline*}
\frac{1}{L}\int_{M}^{M+L}\chi_{V^1_{\epsilon}(B_v)
\times
V^2_\epsilon(C_v)}(\phi_{t+V+v}x,\psi_{t+V}y)\,dt \\ =\frac{1}{L}\int_{M+V}^{M+V+L}\chi_{\phi_{-v}(V^1_{\epsilon}(B_v))
\times
V^2_\epsilon(C_v)}(\phi_{t}x,\psi_{t}y)\,dt,
\end{multline*}
hence the estimate in~\eqref{eq:suf} follows from~\eqref{eeq3} with $M=M+V$ and $L=L$.

By a similar reasoning, we get
\begin{equation}\label{rhot}\rho(\phi_{-v}(B_v)\times C_v)<\rho(B_v\times C_v)+\frac{\eta}{2},
\end{equation}
so putting together~\eqref{disjon} and\eqref{rhot} we derive the estimate
\[|\rho(\phi_{-v}(B_v)\times C_v)-\rho(B_v\times C_v)|<\frac{\eta}{2}.\]
This however contradicts~\eqref{iddisT}, hence the argument is complete.
\end{proof}

\begin{thebibliography}{99}
\bibitem[16]{Gla} E. Glasner, \emph{Ergodic theory via joinings}, Mathematical Surveys and Monographs 101, American Mathematical Society, Providence, RI, 2003.
\bibitem[22]{KLU} A. Kanigowski, M. Lema\'nczyk and C. Ulcigrai, \emph{On disjointness properties of some parabolic flows}, 2018, \url{https://arxiv.org/abs/1810.11576}.
\bibitem[15]{Ga-Le-Sch} P. Gabriel, M. Lema\'nczyk and K. Schmidt, Extensions of cocycles for hyperfinite actions and applications, \emph{Monatshefte f\"ur Mathematik} 123 (1997), no. 3, 209--228, DOI 10.1007/BF01318233.
\end{thebibliography}
\end{document}
